\documentclass[a4paper,12pt]{article}
\usepackage[utf8]{inputenc}
\usepackage[T5]{fontenc}
\usepackage[vietnamese]{babel}
\usepackage{amsmath}
\usepackage{amsfonts}
\usepackage{amssymb}
\usepackage{graphicx}
\usepackage{hyperref}
\usepackage{geometry}

% Page layout
\geometry{left=2.5cm, right=2.5cm, top=2.5cm, bottom=2.5cm}

\title{\textbf{HƯỚNG DẪN SỬ DỤNG VÀ THUYẾT MINH \\ ỨNG DỤNG WEB MINH HỌA SẮP XẾP NGOẠI}}
\author{Học viên: [Tên Học Viên]}
\date{\today}

\begin{document}

\maketitle

\begin{abstract}
Tài liệu này cung cấp hướng dẫn cài đặt và sử dụng Ứng dụng Web minh họa giải thuật Sắp xếp Ngoại (External Sort). Ứng dụng chạy trực tiếp trên trình duyệt, mô phỏng quá trình phân mảnh và trộn K-đường (K-Way Merge Sort) để xử lý lượng dữ liệu vượt giới hạn bộ nhớ (RAM).
\end{abstract}

\section{Yêu cầu Hệ thống}
\begin{itemize}
    \item Trình duyệt web hiện đại hỗ trợ HTML5/ES6 (Google Chrome, Microsoft Edge, Safari, Firefox).
    \item Không yêu cầu cài đặt phần mềm bên thứ ba hay môi trường lập trình (như Python, Node.js).
\end{itemize}

\section{Hướng dẫn Sử dụng}
\begin{enumerate}
    \item \textbf{Khởi chạy ứng dụng:} Mở tệp \texttt{web\_app/index.html} bằng trình duyệt web.
    \item \textbf{Chuẩn bị Dữ liệu:}
          \begin{itemize}
              \item Tại mục "Hoặc sinh file mẫu", nhập \textbf{Số phần tử (N)} và nhấn \textbf{Tạo \& Tải Về Mẫu}. 
              \item Hệ thống sẽ tự động tạo ra một tệp nhị phân \texttt{sample\_N.bin} (mỗi phần tử định dạng Float64 8-bytes).
          \end{itemize}
    \item \textbf{Nạp Dữ liệu:} Nhấn nút chọn tệp và tải lên tệp \texttt{.bin} vừa tạo.
    \item \textbf{Cấu hình Bộ nhớ:} Tại trường \textbf{Kích thước RAM}, nhập giới hạn số phần tử tối đa mà bộ nhớ ảo có thể chứa (ví dụ: 5).
    \item \textbf{Vận hành:} Nhấn nút \textbf{Bắt đầu} để quan sát tiến trình. Người dùng sử dụng các nút \textbf{Tạm Dừng / Bước Kế} nhằm điều khiển tốc độ mô phỏng theo từng lượt thao tác.
\end{enumerate}

\section{Diễn giải Luồng Thuật toán}
Nhờ vào chức năng giả lập, người dùng có thể trực tiếp quan sát thuật toán được thi hành theo hai phân đoạn chính:

\subsection{Giai đoạn 1: Phân mảnh dữ liệu (Split \& Local Sort)}
Thuật toán chia tệp nhị phân gốc thành các khối (chunk) có biểu đồ kích thước bằng với bộ nhớ RAM. Lần lượt từng khối được đọc vào RAM (các khối được highlight viền xanh), sắp xếp cục bộ tại RAM, và xuất ra hệ thống tập tin tạm dưới dạng \texttt{chunk\_i.bin}. 

\subsection{Giai đoạn 2: Trộn K-đường (K-Way Merge)}
Bộ nhớ RAM thực hiện khởi tạo cấu trúc Min-Heap bằng cách nạp 1 phần tử nhỏ nhất từ mỗi tệp tạm. Thuật toán tiến hành tìm giá trị nhỏ nhất trong thể tích RAM hiện hành (highlight vàng vòng lặp so sánh). Số nhỏ nhất được chuyển xuống mảng kết quả (Output Disk), sau đó RAM lập tức kéo phần tử tiếp theo từ chính tệp tạm tương ứng vào để lấp đầy ô trống. Tiến trình lặp lại cho đến khi toàn bộ số cạn kiệt, cho ra một tệp nhị phân kết xuất hoàn chỉnh.

\end{document}
